\section{The Propers: Textual and Liturgical Commentary}
\runninghead{Textual and Liturgical Commentary}

\subsection{Introit (\textit{Ps.~24:1--3, 4b})}

Psalm~24 (Heb.~25) is an alphabetical acrostic prayer ascribed in its heading to David, asking God for deliverance from mocking adversaries, guidance in His ways, and pardon for sin (for its attribution and chronological bounds, see \hyperref[sec:date-location]{\work{Scriptural Date and Location}}). In four places the Introit departs from the Gallican Vulgate of Psalm~24:1--4. It omits the psalm title (\latin{In finem. Psalmus David.}, ``Unto the end, a psalm for David'') and the vocative \latin{Domine} (``O Lord'') in verse~1, beginning directly with \latin{Ad te levavi animam meam} (``To thee have I lifted up my soul''). In verse~3b, it retains the older Roman Psalter phrase \latin{qui te exspectant} (``who wait on thee'') where the Gallican Vulgate reads \latin{qui sustinent te} (``who endure'' or ``wait for thee''). It omits the imprecatory half-verse~4a (\latin{Confundantur omnes iniqua agentes supervacue}, ``Let all them be confounded that act unjust things without cause''). In the psalm verse, it sings the Gallican wording \latin{Vias tuas, Domine, demonstra mihi} (``Show me thy ways, O Lord''), followed by the \latin{Gloria Patri} (``Glory be to the Father''), which is retained throughout Advent (RGMR~427--429), and the repetition of the antiphon.

Medieval and modern liturgical commentators linger over this opening chant. John Beleth (\work{Summa de ecclesiasticis officiis}~LV, PL~202, col.~62) notes that the Church begins her liturgical year with \latin{Ad te levavi} as one looking from afar toward the Nativity; William Durandus (\work{Rationale divinorum officiorum}~VI.3, fol.~261r) observes that the cantor raises the melody from a lower pitch to a higher one to rouse sleepers, matching the Epistle's call to wake from sleep. At the stational Basilica of St.~Mary Major, Bl.~Ildefonso Schuster (\work{The Sacramentary}~I, p.~321) hears in \latin{Ad te levavi} the voice of humanity ``cast down, yet full of hope,'' begging the Savior to bring it back along the way of truth and justice to Bethlehem.

\subsection{Collect (\textit{Excita, quaesumus, Domine})}

The Collect \latin{Excita, quaesumus, Domine, potentiam tuam, et veni} (``Stir up, we beseech Thee, O Lord, Thy power, and come'') stands at the head of Advent in the Gregorian Sacramentary (\work{Hadrianum}, Wilson~1915, p.~113, \latin{Dominica prima}) with the same ending, \latin{Qui vivis et regnas cum Deo Patre} (``Who livest and reignest with God the Father''). By contrast, the older Gelasian Sacramentary arranged five pre-Gregorian Advent Masses without this opening prayer sequence (\work{Sacramentarium Gelasianum}, Wilson~1894, II.LXXX--LXXXIV, pp.~214--218, and p.~lxxi). Unlike Roman Collects addressed to the Father through the Son, its \latin{Qui vivis} conclusion addresses the prayer directly to Christ Himself; as Bl.~Ildefonso Schuster (\work{The Sacramentary}~I, p.~321) notes, the Church asks our Lord to come to her aid and save her by His power from the dangers to which we are exposed through sin.

Its opening imperative echoes Psalm~79:3 (Heb.~80:3), \latin{Excita potentiam tuam, et veni, ut salvos facias nos} (``Stir up thy might, and come to save us''), without reproducing the verse verbatim. Where the psalm appeals against external foes, the Roman prayer names an inward danger: \latin{ab imminentibus peccatorum nostrorum periculis}, ``from the approaching dangers of our sins.'' Two balanced participial clauses then attribute rescue entirely to Christ: \latin{te mereamur protegente eripi, te liberante salvari}---that under His protection we may deserve to be snatched away, and by His freeing power be saved. Under the universal 1962 rubrics, this single Collect stands alone, the pre-1955 seasonal prayers formerly added after it having been abolished in 1960 (RGMR~434, 454, 457~d; \work{Variationes}~II.13, IV.18).

\subsection{Epistle (\textit{Rom.~13:11--14a})}

In St.~Paul's Letter to the Romans (for its setting and date, see \hyperref[sec:date-location]{\work{Scriptural Date and Location}}), this paragraph follows directly on his teaching that Christians owe no debt except mutual love, because love fulfills the Law (Rom.~13:8--10). In the 1962 Missal, the Epistle prepends the address \latin{Fratres:} (``Brethren:''), omits the Vulgate words \latin{Et hoc\dots{} tempus} (``And that, knowing the season'') from verse~11, and ends at \latin{sed induimini Dominum Jesum Christum} (``but put ye on the Lord Jesus Christ''), omitting verse~14b (\latin{et carnis curam ne feceritis in desideriis}, ``and make not provision for the flesh in its concupiscences''). The reading therefore closes directly on Christ's name, whereas St.~Augustine's conversion in \work{Confessions}~VIII.12.29 turned on reading through the end of verse~14.

The liturgical appointment of \latin{Scientes quia hora est} (``Knowing that it is now the hour'') in Advent reaches back at least to the Carolingian period: it is recorded by Amalarius of Metz (\work{Liber de ordine antiphonarii}~XXX; censured by a synod at Kiersy [Quierzy] in 838 for his opinion on the signification of the parts of the divided Host at Mass), and later expounded by Rupert of Deutz (\work{De divinis officiis}~III.2, PL~170, col.~58) and William Durandus (\work{Rationale}~VI.3, fol.~261r).

\subsection{Gradual (\textit{Ps.~24:3b, 4b})}

Following the Epistle, the Gradual returns to Psalm~24:3b and~4b. In the respond, the Missal drops the introductory conjunction \latin{etenim} (``for''), keeps Old Latin \latin{qui te exspectant} (``who wait on thee''), and adds the closing vocative \latin{Domine} (``O Lord''). The Introit verse and the Gradual verse, sung minutes apart in the same Mass, use two different Latin verbs for Psalm~24:4b: where the Introit verse sings the Gallican Vulgate \latin{Vias tuas, Domine, demonstra mihi} (``Show me thy ways, O Lord''), the Gradual keeps the older Roman Psalter wording \latin{Vias tuas, Domine, notas fac mihi} (``Make thy ways known to me''), retaining \latin{et semitas tuas edoce me} (``and teach me thy paths'') in both places. The liturgical repertoire thus preserves two Latin strata side by side in a single Mass.

\subsection{Alleluia (\textit{Ps.~84:8})}

Psalm~84 (Heb.~85), inscribed to the Sons of Korah, recalls God's past forgiveness of His people's iniquity and restoration of Jacob's captivity (vv.~2--4) before pleading for complete renewal and peace (vv.~5--14; see \hyperref[sec:date-location]{\work{Scriptural Date and Location}}). From the center of that plea, the Sunday Mass takes verse~8: \latin{Ostende nobis, Domine, misericordiam tuam: et salutare tuum da nobis} (``Show us, O Lord, thy mercy; and grant us thy salvation''). By the printed 1962 rubric standing directly before this chant (p.~2), the Alleluia and its verse are sung on Sunday, but whenever this Sunday Mass is resumed on an Advent weekday (RGMR~299), the Alleluia is omitted and only the Gradual is said, without replacing it by a Tract.

Why does the Roman liturgy sing \latin{Alleluia} on Advent Sundays while keeping \latin{Gloria in excelsis} silent? William Durandus (\work{Rationale}~VI.2, fol.~256v) puts the question directly: if Advent is a fast of mourning, every song of rejoicing should cease; if of rejoicing, the \latin{Gloria} too should be sung. His solution is that this fast is ``partly of rejoicing and partly of mourning'' (\latin{partim exultationis, et partim moeroris}), just as John Beleth (\work{Summa}~LXIV, PL~202, col.~71) explains that \latin{Alleluia} is not silenced as in Septuagesima because Advent is ``partly of gladness'' (\latin{partim laetitiae}).

\subsection{Gospel (\textit{Luke~21:25--33})}

In St.~Luke's Gospel, chapter~21 unfolds during Passion Week as Jesus teaches daily in the Temple courts and spends the nights on the Mount of Olives (Luke~21:5, 37--38; see \hyperref[sec:date-location]{\work{Scriptural Date and Location}}). Admiring onlookers prompt His prophecy of the Temple's ruin, His warnings against deceivers and persecution, and His foretelling of Jerusalem's siege (vv.~6--24). At verse~25, the Sunday Gospel opens with the liturgical opening \latin{In illo tempore: Dixit Iesus discipulis suis:} (``At that time, Jesus said to his disciples:''), replacing the Vulgate's bare conjunction \latin{Et} (``And''). The reading runs through verse~33, spanning celestial signs, the distress of nations at the roaring sea, men withering with fear, the shaking of the heavenly powers, the Son of Man coming in a cloud with power and majesty, the command to look up and lift up your heads, the parable of the fig tree and all the trees, and Christ's pledge that His words will outlast heaven and earth. (At verse~29, the 1899 American Douay--Rheims edition prints ``And he spoke to them in a similitude,'' whereas the Challoner Gutenberg 1581 edition omits the preposition and reads ``a similitude.'')

Two details distinguish St.~Luke's wording from the parallel discourse in Matthew~24. First, Luke speaks specifically of the roaring of the sea and waves (\latin{sonitus maris, et fluctuum}) and of men withering away (\latin{arescentibus hominibus}) for fear and expectation. Second, in verse~29 Luke alone adds ``and all the trees'' (\latin{et omnes arbores}) beside the fig tree. William Durandus (\work{Rationale}~VI.3, fol.~261v) pairs this Gospel's promise that ``the kingdom of God is at hand'' (\latin{prope est regnum Dei}) directly with the Epistle's ``our salvation is nearer'' (\latin{propior est nostra salus}).

\subsection{Offertory (\textit{Ps.~24:1--3})}

As bread and wine are placed on the altar, the choir repeats word for word the entire Introit antiphon: \latin{Ad te levavi animam meam\dots{} etenim universi, qui te exspectant, non confundentur}. Rupert of Deutz (\work{De divinis officiis}~III.2, PL~170, col.~58) points out this return to the entrance antiphon, reading the repetition as a sign that the Church continues to cry for God's guidance until she reaches the heavenly Jerusalem. Bl.~Ildefonso Schuster (\work{The Sacramentary}~I, p.~322) calls Psalm~24 ``the typical psalm of the first Sunday of Advent,'' noting its closing assurance that none who wait upon the Lord in watchful trust will be confounded.

\subsection{Secret (\textit{Haec sacra nos, Domine})}

Preserved verbatim under \latin{Super oblata} (``Over the offerings'') in the Gregorian Sacramentary (\work{Hadrianum}, Wilson~1915, p.~113), the Secret states its petition in a single Latin sentence: \latin{Haec sacra nos, Domine, potenti virtute mundatos, ad suum faciant puriores venire principium} (``Cleansing us by their mighty power, may these Holy Mysteries, O Lord, make us come more pure before Thee who art their author''). The priest says it silently except for the closing \latin{Per omnia saecula saeculorum} (``World without end,'' RGMR~480--481). Bl.~Ildefonso Schuster (\work{The Sacramentary}~I, p.~322) hears in its concise wording ``the ancient idea that runs through the Roman epiklesis,'' beseeching the Savior that the power of the Eucharistic sacrifice purify our minds so we may worthily receive Him who is its Author (\latin{principium}).

\subsection{Communion (\textit{Ps.~84:13})}

At Communion, the Mass returns to Psalm~84 and sings verse~13 without the Vulgate's initial \latin{Etenim} (``For''): \latin{Dominus dabit benignitatem: et terra nostra dabit fructum suum} (``The Lord will give goodness: and our earth shall yield her fruit''). Where the Gallican Vulgate and Roman Missal sing \latin{benignitatem} (``goodness''), St.~Augustine's Old Latin Psalter read \latin{suavitatem} (``sweetness''). Commentators read the verse along four distinct lines: St.~Robert Bellarmine and Sicard of Cremona (\work{Mitrale}~V.2, PL~213, col.~205) identify \latin{benignitas} with the Holy Spirit and \latin{terra nostra} with the Virgin Mary yielding Christ; Dom Prosper Gu\'eranger (\work{The Liturgical Year}~I, p.~137) likewise calls Mary ``our Earth \dots\ made fruitful by the dew of heaven'' after Isaiah while connecting the antiphon to Christ giving Himself by anticipation in Communion; St.~Augustine, reading \latin{suavitatem}, expounds the verse as interior delight in righteousness poured into hearts ploughed by confession; Bl.~Ildefonso Schuster (\work{The Sacramentary}~I, pp.~322--323) reads it as thanksgiving for the Eucharistic pledge refreshing hearts made arid by sin; and William Durandus (\work{Rationale}~VI.3, fol.~261v) takes \latin{benignitatem} as good affection (\latin{affectum bonum}) and \latin{terra nostra} as our flesh yielding good works.

\subsection{Postcommunion (\textit{Suscipiamus, Domine, misericordiam tuam})}

Attested under \latin{Ad complendum} (``At the completion'') in the Gregorian Sacramentary (\work{Hadrianum}, Wilson~1915, p.~113), the Postcommunion differs from the earliest Gregorian manuscript in one word: the \work{Hadrianum} reads \latin{et reparationis} (``and\dots{} of our reparation''), whereas the Tridentine and 1962 Missals print the purpose clause \latin{ut reparationis} (``so that\dots''). Its opening words echo Psalm~47:10 (Heb.~48:10; modern English~48:9), \latin{Suscepimus, Deus, misericordiam tuam in medio templi tui} (``We have received thy mercy, O God, in the midst of thy temple''), turning the psalm's past tense into the petition \latin{Suscipiamus} (``May we receive''). Even after sacramental Communion, the Church prays to receive God's mercy within His temple so as to go before ``the coming solemnities of our reparation'' (\latin{reparationis nostrae ventura solemnia}) with fitting honor. William Durandus (\work{Rationale}~VI.3, fol.~259r) cites this Postcommunion explicitly as marking the preparation of the Savior's ways, and Bl.~Ildefonso Schuster (\work{The Sacramentary}~I, p.~323) calls it ``the true \latin{Eucharistia} in its etymological sense,'' inspired by Psalm~47, asking that the divine virtue of the heavenly bread prepare the assembly suitably for ``the approaching feast of our common ransom.''

\subsection{The Governed Ordinary of the Sunday}

Five variable rules of the 1962 Ordinary accompany these ten proper texts. First, violet vestments and the silence of \latin{Gloria in excelsis} throughout Advent \latin{de Tempore} (1962 printed rubric, p.~1; RGMR~431--432) mark penitential expectation. Medieval commentators give complementary reasons for silencing the hymn: William Durandus (\work{Rationale}~VI.2, fol.~256r) writes that the Nativity hymn is withheld so it may be sung more devoutly and eagerly on Christmas night; Amalarius of Metz (\work{Liber officialis}~III.40, PL~105, col.~1159; censured by a synod at Kiersy [Quierzy] in 838 for his opinion on the signification of the parts of the divided Host at Mass) records its Advent omission as signifying the time of the Old Testament, an explanation echoed by Honorius Augustodunensis (\work{Gemma Animae}~III.1, PL~172, col.~642, recalling the righteous waiting before Christ's coming) and Rupert of Deutz (\work{De divinis officiis}~III.2, PL~170, col.~58, pointing to the age that still awaited the joy of peace as well as the Church's waiting for final salvation); John Beleth (\work{Summa}~LXIV, PL~202, col.~71) explains differently that \latin{Gloria in excelsis} (belonging to the Nativity) and \latin{Ite, missa est} (belonging to the Resurrection) are both omitted in Advent because both mysteries are then eagerly awaited (\latin{quod tunc utrumque avide exspectetur}). Second, while medieval writers such as Rupert (III.2) and Beleth (LXIV) also note the omission of \latin{Ite, missa est} (``Go, the Mass is ended'') whenever \latin{Gloria} was omitted, the 1962 Missal retains \latin{Ite, missa est} on Advent Sundays, restricting \latin{Benedicamus Domino} (``Let us bless the Lord'') to the Holy Thursday evening Mass followed by solemn reposition and to other Masses followed by a procession (RGMR~507).

Third, the Sunday \latin{Credo} is required after the Gospel (printed rubric, p.~2; RGMR~475--476), confessing both the Son's incarnation by the Holy Spirit of the Virgin Mary and His return in glory to judge the living and the dead. Fourth, because the 1962 Missal contains no proper Preface of Advent, its printed rubric on page~2 appoints the Preface of the Most Holy Trinity for all Sundays of Advent (n.~1082, printed p.~293; RGMR~494~b), reserving the Common Preface for Advent weekdays resuming the Sunday Mass (RGMR~482, 498). Confessing the Father together with His only-begotten Son and the Holy Spirit---``distinctness in persons, oneness in essence, and equality in majesty''---the Trinity Preface joins the congregation to the unceasing praise of angels, archangels, cherubim, and seraphim. Fifth, the ordinary Roman Canon follows without seasonal proper clauses (RGMR~500--501), and the ordinary Last Gospel, John~1:1--14 (RGMR~509--510), concludes Mass with the eternal Word made flesh dwelling among us.
